\documentclass[11pt]{article}
\usepackage[margin=1in]{geometry}
\usepackage{amsmath,amssymb,amsthm}
\usepackage{boldtensors}
\usepackage{hyperref}
\usepackage{booktabs}

\newcommand{\dd}{\mathrm{d}}

\title{Design: Mortar Transfer for Impedance-Overlap Schwarz \\
on Non-Conforming Meshes}
\author{Alejandro Mota}

\begin{document}
\maketitle

\section{Problem statement}

The impedance-overlap transmission condition
\begin{equation}
  ~t_k - ~t_j
  + ~Z\,\bigl(\dot{~u}_k - \dot{~u}_j\bigr)
  + \alpha\,\bigl(~u_k - ~u_j\bigr) = ~0
  \qquad \text{on } \Gamma_k
  \label{eq:tc}
\end{equation}
is now exactly conservative on conforming meshes (0.08\% PU-energy error
over 1\,ms with the consistent partner traction, matching the
exact-DBC reference).  On the non-conforming cantilever pair
($h = 8.47$ vs.\ $6.35$\,mm, 4:3) the net error also looks small
($\pm 1$\%), but the interface-work instrumentation shows that this is
a cancellation of two error channels, each of which moves roughly
\emph{half the system energy per millisecond}:

\begin{center}
\small
\begin{tabular}{lccc}
\toprule
1\,ms, II, PU metric & conforming & non-conforming $h$ & non-conforming $h/2$ \\
\midrule
dashpot (spurious) work & $-42.6$\,J \,(2.8\%) & $-606.5$\,J \,(40\%) & $-598.5$\,J \,(40\%) \\
traction-channel work & --- & $+677$\,J \,(45\%) & $+768$\,J \,(51\%) \\
net PU-energy error & $1.5$\% $\to$ 0.08\% & $-0.75$\% (growth) & $+1.5$\% (loss) \\
mean velocity jump (RMS) & 8.9\,m/s & 40\,m/s (growing 15$\to$60) & 27\,m/s \\
\bottomrule
\end{tabular}
\end{center}

Three facts drive the design.  (i) The channels do \emph{not} shrink
under mesh refinement (dashpot work unchanged at $h/2$), so resolution
cannot fix the interface; the pointwise-interpolation transfer feeds
the condition an $O(1)$ error on mesh-scale content.  (ii) The net
error is the difference of two large channels and even flips sign
between $h$ and $h/2$ --- it is not robust, which is the production
fragility.  (iii) The velocity jumps grow through the run: each
interface pass leaves mesh-scale residue that reverberates, a bounded
version of the same feedback that makes DBC overlap coupling
exponentially unstable.

The theoretical requirements are Gander--Halpern--Nataf's three
conditions for discrete energy convergence of waveform-relaxation
coupling on nonmatching grids: (i) Robin (not Dirichlet) transmission,
(ii) consistent boundary-cell discretization of the transmission
operator, (iii) an $L^2$-contractive intergrid transfer.  The
impedance-overlap BC satisfies (i); its weak (surface-mass-weighted)
assembly satisfies (ii); pointwise interpolation violates (iii).
Replacing the transfer with the $L^2$ (mortar) projection onto the
receiving trace space satisfies all three simultaneously.  Note that
weak-$L^2$ enforcement of \emph{Dirichlet} overlap coupling did not
remove the DBC instability --- consistent with the theory, since DBC
violates (i) regardless of the transfer; the combination proposed here
is the first that meets all three conditions.

\section{Operator inventory (existing machinery)}

All operators needed for the first stage already exist:

\begin{center}
\small
\begin{tabular}{lll}
\toprule
operator & code & status \\
\midrule
$~W_k = \int_{\Gamma_k} ~N_k ~N_k^\top \dd\Gamma$ &
  \texttt{get\_square\_projection\_matrix} &
  already on the BC \\[0.3em]
$~L_{kj} = \int_{\Gamma_k} ~N_k(~x)\, ~N_j^{\mathrm{vol}}(~x)^\top \dd\Gamma$ &
  \texttt{get\_overlap\_rectangular\_projection\_matrix} &
  exists (weak-DBC path); \\
  & & needs signature relaxed \\[0.3em]
$~P_{kj} = ~W_k^{-1} ~L_{kj}$ &
  weak-DBC \texttt{dirichlet\_projector} &
  same construction \\[0.3em]
surface-to-surface $~R_{kj}$ &
  \texttt{get\_rectangular\_projection\_matrix} &
  exists (nonoverlap path); \\
  & (\texttt{project\_point\_to\_side\_set}) & needed in stage 2 \\[0.3em]
partner one-sided weak flux &
  \texttt{build\_consistent\_traction\_patch!}, &
  exists (conforming); \\
  & \texttt{consistent\_partner\_traction} & generalized in stage 2 \\
\bottomrule
\end{tabular}
\end{center}

$~L_{kj}$ is assembled by quadrature over $\Omega_k$'s boundary facets;
at each quadrature point the partner \emph{volume} element is located
(\texttt{find\_point\_in\_mesh}) and its shape functions evaluated ---
exactly the mortar cross-mass for an overlap boundary, which is
interior to the partner and therefore has no partner surface mesh.

Two structural properties come for free.  First,
$~P_{kj} = ~W_k^{-1} ~L_{kj}$ is the $L^2(\Gamma_k)$-orthogonal
(Galerkin) projection onto $\Omega_k$'s trace space, hence
non-expansive in $L^2$: GHN condition (iii) holds by construction,
for any quadrature, because it is a projection.  Second, partition of
unity of the partner shape functions gives
$~L_{kj} \mathbf{1} = ~W_k \mathbf{1}$ for \emph{any} quadrature rule,
so constants (rigid translations) transfer exactly regardless of
quadrature accuracy.

\section{Design}

\subsection{Stage 1: mortar projection of the transferred fields}

Replace pointwise interpolation of the three partner fields with the
mortar projection, leaving everything downstream of the nodal values
unchanged:
\begin{equation}
  \hat{\dot{~u}}_j = ~P_{kj}\, \dot{~u}_j, \qquad
  \hat{~u}_j = ~P_{kj}\, ~u_j, \qquad
  \hat{~\sigma}_j = ~P_{kj}\, ~\sigma_j^{\mathrm{rec}}
  \quad \text{(six Voigt components),}
  \label{eq:stage1}
\end{equation}
then form the traction $\hat{~t}_j = \hat{~\sigma}_j\, ~n_k$ and the
tensor-impedance term nodally exactly as the current code does.  This
is a drop-in change: \texttt{apply\_bc\_detail}'s per-node
interpolation loop becomes three (or nine, per component) small dense
matrix--vector products with a projector precomputed at setup.  The
P/S-split tensor $~Z(~n)$ continues to act at the receiving nodes, so
no re-derivation of the self terms or the tangent is needed; the
variational crime of nodal application is unchanged from the current
scheme and identical on both branches.

Implementation touch points:
\begin{itemize}
\item \textbf{YAML}: new BC option \texttt{transfer:} with values
      \texttt{pointwise} (current behavior, initial default) and
      \texttt{mortar}.
\item \textbf{Types}: field \texttt{mortar\_projector::Matrix\{Float64\}}
      ($n_k \times N_j$, dense; boundary node count times partner node
      count, small for realistic interfaces) on
      \texttt{SolidMechanicsImpedanceOverlapSchwarzBoundaryCondition}.
\item \textbf{Setup} (\texttt{compute\_impedance\_overlap\_schwarz\_projectors!}):
      relax \texttt{get\_overlap\_rectangular\_projection\_matrix}'s
      signature to accept the impedance BC (it only touches fields both
      BC types share), build $~L_{kj}$, factorize the existing
      $~W_k$ (Cholesky) once, store $~P_{kj}$.  Rebuilt on BC swaps,
      like the consistent-traction patch.
\item \textbf{Evaluation}: factor the partner-field sampling shared by
      \texttt{apply\_bc\_detail} and
      \texttt{report\_impedance\_interface\_work} into one helper that
      dispatches on the transfer mode, so the instrumentation always
      measures what the BC actually applies.
\item \textbf{Precedence}: on node-aligned interfaces $~L_{kj}$ has a
      single unit entry per row and $~P_{kj}$ reduces to the identity;
      the consistent-traction patch (\texttt{partner traction: auto})
      continues to take priority for the traction, and the kinematic
      projection becomes a no-op.  The conforming limit is therefore
      automatic.
\end{itemize}

Quadrature note: the integrand of $~L_{kj}$ is only piecewise smooth
on $\Gamma_k$ (partner elements change within a receiving facet), so
the default facet rule carries a consistency error.  Because
contractivity and constant-transfer hold for any rule, this error
degrades accuracy, not stability.  The cheap knob is an elevated or
subdivided facet rule for the $~L_{kj}$ assembly only (setup cost);
a true common-refinement (2D polygon overlay on the flat interface)
is the exact alternative and is deferred until measurements justify
it.

What stage 1 answers: how much of the 600\,J churn is the
\emph{transfer} (interpolation aliasing of mesh-scale content) versus
the \emph{recovered-stress traction}.  The conforming ladder bounds
what recovery-quality traction can achieve ($\approx$1.5\%-class); if
mortar transfer restores $O(h)$ scaling of the channels, refinement
plus stage 2 closes the rest.

\subsection{Stage 2: consistent traction across the interface}

The conforming analysis showed the traction channel is the dominant
residual once transfer is clean, and that nodal recovery cannot
represent the traction carried by mesh-scale content.  The consistent
(one-sided partial assembly) traction requires a partner facet surface
bounding a one-sided element patch; on a non-conforming interface
$\Gamma_k$ cuts through partner elements, so the patch does not exist
on $\Gamma_k$ itself.  Two generalizations, in order of increasing
cost:

\paragraph{2a: offset facet surface + surface mortar.}  Take
$\tilde\Gamma_j$ = the partner facet surface immediately exterior to
$\Gamma_k$ (the boundary between the first exterior element layer and
the rest; the patch-construction logic already classifies these
elements).  The consistent flux is exact on $\tilde\Gamma_j$ by the
same partial assembly as the conforming case, tested with partner
surface functions.  Transfer the resulting weak traction to
$\Gamma_k$'s trace space with the surface-to-surface mortar
($\texttt{get\_rectangular\_projection\_matrix}$ generalized to accept
the internal surface in place of a side set, or an equivalent
assembled operator), i.e.\
$~W_k \hat{~t}_j = ~R_{k\tilde\jmath}\, ~W_{\tilde\jmath}^{-1}
\lambda_{\tilde\jmath}$.
The surface offset introduces an $O(h_j)$ modeling error
($\tilde\Gamma_j$ is within one partner element of $\Gamma_k$), but it
acts on the \emph{exact discrete flux} rather than on recovered
stress; whether that trade wins is an experiment, not a theorem ---
measure before committing further.

\paragraph{2b: cut-cell partial assembly (flat interfaces).}  For a
planar $\Gamma_k$ (the common production configuration), integrate the
partner's momentum residual over the plane-clipped parts of the cut
elements, testing with $\Omega_k$'s surface functions extended
constant along the interface normal.  This reproduces the
conforming construction exactly in the limit and removes the offset
error of 2a, at the cost of new machinery (hexahedron--half-space
clipping and quadrature on the clipped polyhedra).  Deferred unless 2a
measures short of the target.

\subsection{Validation and targets}

The metric is the instrumentation's channel decomposition, not the net
energy: success is driving the churn itself down, since the current
scheme already fakes a small net by cancellation.

\begin{enumerate}
\item \textbf{Conforming regression}: mortar transfer on the
      node-aligned pair must reproduce the pointwise results exactly
      ($~P = ~I$); the consistent-traction 0.08\% result must be
      unaffected.  Extend test 103/104 asserts with the projector
      identity check.
\item \textbf{Non-conforming, stage 1}: dashpot work
      $\ll 600$\,J at $h$, and restored $O(h)$ decay at $h/2$;
      velocity-jump growth curve flattened.  Expected end state
      $\approx$ conforming recovered-stress class (few percent), now
      robust rather than cancellation-balanced.
\item \textbf{Non-conforming, stage 2}: traction channel reduced
      toward the consistent-traction class; net PU error
      $\lesssim 1$\% with both channels individually small; sign and
      magnitude stable under refinement.
\item \textbf{Stability margin}: rerun the impedance-scale sweep ---
      with contractive transfer the growth branch above scale 1 should
      soften, and the DBC-limit blow-up should set in later; this is
      the direct measure of production robustness.
\item \textbf{Cost}: setup $O(n_{\mathrm{facets}} \times n_{\mathrm{qp}})$
      point locations (same as the weak-DBC path); per-iteration cost
      \emph{decreases} (dense mat-vec replaces per-node interpolation
      loops; no recovery solve once stage 2 lands).
\end{enumerate}

\section{Summary}

Stage 1 is a low-risk, high-information change that reuses the square
($~W_k$) and rectangular ($~L_{kj}$) projection machinery verbatim:
mortar-project $\dot{~u}_j$, $~u_j$, and $~\sigma_j^{\mathrm{rec}}$
onto the receiving trace space and keep the rest of the impedance
pipeline unchanged.  It makes the transfer $L^2$-contractive by
construction, completing GHN's three conditions for the first time in
this code, and its measurements decide between the two stage-2
traction designs.  The conforming limit and the consistent-traction
path are preserved automatically.

\end{document}
